%% 第四章
\chapter{能力优化方法与实验}
\chaptermark{能力优化方法与实验}

\section{UAV-DualCog-Train训练集构建}

\subsection{构建目标与场景隔离}

第3章的评测结果表明，当前多模态大模型在 UAV-DualCog 上的主要问题是难以同时保持答案正确、证据可靠、跨媒介稳定以及自我状态认知与环境情况认知之间的协调发展。尤其是在图像任务中，模型常表现为选项判断相对较强、空间证据明显偏弱；在双认知维度上，自我状态认知长期落后于环境情况认知。因此，本文在测试基准之外进一步构建了面向后训练优化的图像任务训练集 UAV-DualCog-Train，并据此设计“监督微调+联合 GRPO”的能力优化方法。

UAV-DualCog-Train 使用与测试集严格隔离的训练场景构建，包含 \texttt{env\_1}、\texttt{env\_3}、\texttt{env\_5}、\texttt{env\_15}、\texttt{env\_18} 和 \texttt{env\_22} 六个场景，具体训练/测试场景划分见\autoref{tab:app-scene-split}。训练样本覆盖地标参照位置推理、未来观察预测、自我参照位置推理和地标驱动动作决策四类任务，对应自我状态认知与环境情况认知两条能力主线。当前优化阶段优先从图像任务起步，主要基于两点考虑：其一，第3章已经表明图像侧最集中地暴露了“语义判断与空间证据脱节”的问题；其二，图像任务适合将“离散答案+目标框证据”统一写成结构化监督目标，为后续监督微调和奖励优化共用同一套输出解析接口提供便利。

% 去除最后一列
\begin{table}[htbp]
    \centering
    \caption{UAV-DualCog-Train-V1 训练场景与样本规模统计}
    \label{tab:train-scene-stats}
    \footnotesize
    \begin{tabularx}{\textwidth}{@{}>{\centering\arraybackslash}p{2.1cm}>{\centering\arraybackslash}p{2.4cm}X@{}}
        \toprule
        场景编号 & 有效地标数 & 训练样本数\\
        \midrule
        \texttt{env\_1} & 80 & 3840\\
        \texttt{env\_3} & 31 & 1488\\
        \texttt{env\_5} & 8 & 384\\
        \texttt{env\_15} & 38 & 1824\\
        \texttt{env\_18} & 54 & 2592\\
        \texttt{env\_22} & 23 & 1104\\
        \midrule
        合计 & 234 & 11232 \\
        \bottomrule
    \end{tabularx}
\end{table}

\subsection{训练记录格式与样本组织}

为保证训练目标与评测协议一致，训练记录与真实多模态输入直接对齐。每条样本会被转换为统一的 \texttt{prompt}、\texttt{completion} 和 \texttt{images} 三元结构，其中 \texttt{prompt} 保留 benchmark 原有的系统提示与用户提示，\texttt{images} 以与评测阶段一致的顺序存放图像路径，\texttt{completion} 则使用结构化 JSON 监督目标：
\begin{verbatim}
{
  "answer_option_id": "B",
  "bbox_xyxy_norm": [0.112193, 0.400801, 0.382753, 0.698569]
}
\end{verbatim}
对于 \texttt{self\_where}、\texttt{env\_where} 和 \texttt{env\_how} 三类任务，训练输入使用“地标参考图 + 查询图像”两图组合；对于 \texttt{self\_what} 任务，训练输入则保留“地标参考图 + 候选图像序列”的多图结构。与此同时，样本记录中还会缓存 \texttt{answer\_bbox\_xyxy\_1000} 等辅助字段，供奖励计算阶段直接复用。这样的设计使训练阶段围绕“结构化答案输出、空间证据可核验”的目标展开。训练数据字段、划分策略和数据准备伪代码见\autoref{app:train-dataset-construction}，训练入口与关键实现片段见\autoref{app:optimization-entry}。

\subsection{训练集规模}

UAV-DualCog-Train-V1 的规模为
\[
234 \times 4 \times 2 \times 6 = 11232,
\]
即 234 个有效地标、4 类任务、2 档难度、每档难度 6 个问答样本。训练集与验证集采用稳定轮转划分策略，在每个 \texttt{(task\_type, difficulty)} 桶内按样本编号排序后，每 10 条样本取 1 条进入验证集，其余进入训练集。相比完全随机划分，这种做法能够在不同任务和难度子集之间保持更稳定的 9:1 比例，同时避免训练集与验证集在样本编号层面发生重叠。

\begin{table}[!htbp]
    \centering
    \caption{UAV-DualCog-Train 数据分布}
    \label{tab:train-dataset-stats}
    \footnotesize
    \setlength{\tabcolsep}{4pt}
    \begin{tabularx}{\textwidth}{@{}>{\raggedright\arraybackslash}X>{\centering\arraybackslash}p{2.2cm}>{\raggedright\arraybackslash}X>{\centering\arraybackslash}p{2.2cm}@{}}
        \toprule
        统计项 & 数值 & 统计项 & 数值 \\
        \midrule
        全量样本数 & 11232 & 训练集样本数 & 10104 \\
        验证集样本数 & 1128 & 训练/验证比例 & 89.96\% : 10.04\% \\
        四类任务样本数 & 每类 2808 & 两档难度样本数 & 各 5616 \\
        双认知分组样本数 & 各 5616 & 每个任务-难度组合 & 1404 \\
        非法目标框数量 & 0 & 参考图路径存在率 & 100\% \\
        查询图路径存在率 & 100\% & 平均提示词长度 & 97.62 词 \\
        平均提示词字符数 & 769.82 & 目标框面积均值 & 0.0645 \\
        \bottomrule
    \end{tabularx}
\end{table}

\autoref{tab:train-dataset-stats} 说明，训练集在任务类型、难度等级和双认知分组上保持严格均衡，不存在单一能力分支或单一难度主导训练的问题。答案选项分布也近似均匀，其中 A、B、C、D 四类选项占比分别为 24.78\%、24.79\%、26.40\% 和 24.03\%，未出现明显类别偏置。目标框几何统计进一步表明，该训练集既包含较小目标，也包含中等尺度和相对显著目标，能够为后续的空间证据学习提供较丰富的监督信号。总体而言，UAV-DualCog-Train-V1 在样本规模、认知平衡、难度平衡和字段完整性四个方面满足后训练优化要求。

\section{能力优化方法}

\subsection{多模态监督微调}

能力优化的第一阶段采用多模态监督微调。该阶段的目标是先为模型建立稳定的任务内输出接口，使其学会围绕 UAV-DualCog 的图像任务生成结构化的“答案+目标框”结果。在训练实现中，图像会由 \texttt{AutoProcessor} 编码为视觉张量参与前向和反向传播。因此，SFT 阶段的核心作用是让模型在多图像输入下稳定输出可解析的结构化答案。

本文选择 \texttt{Qwen/Qwen3.5-4B} 作为当前能力优化的基座模型，主要基于三方面考虑。第一，第3章结果显示 Qwen 3.5-4B 已具备一定图像空间问答基础，其图像总体准确率为 39.0\%。第二，该模型在视频总体准确率和自我状态认知上仍有明显短板，具备清晰的优化空间。第三，4B 规模适合在现有四张 RTX 4090 的算力条件下进行多轮参数高效实验，并便于后续补充消融研究。

在方法上，SFT 阶段使用 \texttt{AutoModelForImageTextToText}、\texttt{AutoProcessor} 和 \texttt{trl.SFTTrainer} 完成多模态 LoRA 微调。LoRA 仅在注意力和前馈网络的关键投影层上注入可训练低秩适配器，使模型在有限显存与训练步数下完成任务适配。训练时采用 \texttt{completion\_only\_loss}，即只对助手端的 JSON 输出部分计算损失，从而把格式学习、答案学习与证据学习统一绑定到同一训练接口中。为兼顾图像细节和显存稳定性，系统在加载训练图像时引入 \texttt{max\_image\_side} 控制项，使用 720 作为长边上限，以保留更充分的空间细节。具体训练入口、配置与关键代码见\autoref{app:optimization-entry}。

% 根据训练集规模和正式训练参数，SFT 阶段的等效全局 batch 为 \(1 \times 4 \times 32 = 128\)，每个 epoch 约为 79 个优化步，总共训练 3 个 epoch，对应约 237 个优化步。SFT 阶段结束后，系统会保存 LoRA 适配器、训练日志、验证指标以及可直接部署的合并模型，为后续联合 GRPO 提供可直接加载的初始化模型。正文仅保留方法与参数口径，具体训练入口、配置与关键代码见\autoref{app:optimization-entry}。

\subsection{联合GRPO优化}

在完成结构化监督学习之后，本文进一步引入联合 GRPO 进行第二阶段优化。该阶段直接在 SFT 检查点基础上，以统一奖励函数同时约束输出格式、离散答案和目标框质量。多模态训练链路已经能够稳定支持真实图像输入和 JSON 输出解析，因此奖励优化可以围绕最终任务目标进行联合校正。

联合 GRPO 的总体奖励形式为
\begin{equation}
R = w_f R_{\text{format}} + w_s R_{\text{mcq}} + w_b R_{\text{bbox}},
\end{equation}
其中 \(R_{\text{format}}\) 衡量模型输出能否被解析为规范 JSON，\(R_{\text{mcq}}\) 衡量离散答案是否正确，\(R_{\text{bbox}}\) 衡量预测目标框与真实目标框的一致性。奖励实现优先解析 \texttt{\{"answer\_option\_id", "bbox\_xyxy\_norm"\}} 形式的 JSON 输出；目标框奖励则基于归一化坐标上的 GIoU 计算，并映射到 \([0,1]\) 区间。该设计与第3章的分析结论对应，即同时缩小“答得对”与“证据也对”之间的差距。

\begin{table}[!htbp]
    \centering
    \caption{能力优化两阶段关键训练参数}
    \label{tab:optimization-training-config}
    \footnotesize
    \setlength{\tabcolsep}{4pt}
    \begin{tabularx}{\textwidth}{@{}p{3.1cm}>{\centering\arraybackslash}X>{\centering\arraybackslash}X@{}}
        \toprule
        参数项 & 多模态 SFT & 联合 GRPO \\
        \midrule
        初始化模型 & Qwen 3.5 4B & SFT 阶段合并模型 \\
        学习率 & \(5\times10^{-5}\) & \(1\times10^{-6}\) \\
        训练长度 & 3 个 epoch，约 237 步 & 600 步 \\
        单卡 batch size & 1 & 1 \\
        梯度累积步数 & 32 & 16 \\
        最大序列/生成长度 & 4096 & 128 \\
        图像长边上限 & 720 & 720 \\
        LoRA 秩 \(r\) / \(\alpha\) & 16 / 32 & 8 / 16 \\
        LoRA dropout & 0.05 & 0.05 \\
        核心优化目标 & 结构化答案学习 & 格式、答案与目标框联合奖励 \\
        \bottomrule
    \end{tabularx}
\end{table}

\autoref{tab:optimization-training-config} 给出了两阶段能力优化的关键参数。联合 GRPO 直接将格式、答案与空间证据纳入同一奖励函数，使训练目标与最终测试基准图像任务更加一致。与此同时，训练端还通过关闭思维展开型模板、限制最大生成长度和保留小批量组相对采样，降低长生成噪声对结构化奖励的干扰。相关奖励实现、伪代码和关键代码片段见\autoref{app:optimization-entry}。

\subsection{训练环境、流程与工程验证}

能力优化实验在一台本地服务器上进行，硬件环境如\autoref{tab:train-hardware}所示。

\begin{table}[htbp]
    \centering
    \caption{能力优化实验训练环境}
    \label{tab:train-hardware}
    \small
    \begin{tabularx}{\textwidth}{@{}>{\centering\arraybackslash}p{3.0cm}>{\centering\arraybackslash}X@{}}
        \toprule
        资源项 & 配置 \\
        \midrule
        CPU & \(2\times\) AMD EPYC 7452，合计 128 vCPUs \\
        内存 & 503 GiB \\
        GPU & \(4\times\) NVIDIA GeForce RTX 4090，单卡 24 GB 显存 \\
        训练框架 & PyTorch、Transformers、PEFT、TRL \\
        训练方式 & 单机多卡多模态 LoRA SFT + 联合 GRPO \\
        \bottomrule
    \end{tabularx}
\end{table}

% 不需要写小样本工程验证，直接写最终的流程即可。
在工程组织上，本文将训练流程统一整理为“数据准备、SFT、Joint GRPO”三个连续阶段，并通过统一调度脚本对输入快照、日志目录、检查点目录和多卡启动方式进行管理。若训练前缺失训练集或验证集 JSONL 文件，系统会自动重新生成数据；若文件已存在，则会将当前使用的数据快照复制到本次运行目录中，确保后续能够准确追溯到对应训练数据版本。每个阶段结束后还会检查是否成功导出适配器与合并模型，避免后续流程建立在不完整中间产物之上。

在正式全量训练之前，本文首先完成了针对新链路的小样本工程验证，确认真实图像加载、视觉编码、SFT 反向传播、GRPO 奖励计算、模型合并和产物落盘均可执行。正式训练完成后，本文进一步基于运行日志和产物目录进行完整性核查，确认 SFT 与 Joint GRPO 两个阶段均达到预设步数，并成功导出对应的适配器与合并模型。小样本验证与完整性核查逻辑见\autoref{app:optimization-entry}。

\begin{table}[!htbp]
    \centering
    \caption{正式训练完成情况与核心过程指标}
    \label{tab:train-run-summary}
    \footnotesize
    \setlength{\tabcolsep}{4pt}
    \begin{tabularx}{\textwidth}{@{}p{1.9cm}>{\centering\arraybackslash}p{1.6cm}>{\centering\arraybackslash}p{1.6cm}>{\centering\arraybackslash}p{2.0cm}X@{}}
        \toprule
        阶段 & 完成步数 & 运行时间/h & 过程指标 & 说明 \\
        \midrule
        SFT & 237/237 & 4.05 & \makecell[c]{train loss\\0.6384} & 完成 3 个 epoch，验证集 \texttt{eval\_loss}=0.5666，图像长边上限为 720。 \\
        Joint GRPO & 600/600 & 10.72 & \makecell[c]{tail-10 reward\\2.4016} & 奖励权重为 \((1.0,1.0,1.0)\)，后段奖励保持稳定，最终导出可评测权重。 \\
        \bottomrule
    \end{tabularx}
\end{table}

\autoref{tab:train-run-summary} 表明，两阶段训练均已完整结束。SFT 阶段使用约 4.05 小时完成 237 个优化步，并在验证集上获得 0.5666 的损失；Joint GRPO 则在约 10.72 小时内完成 600 个优化步，后 10 次记录的平均奖励为 2.4016，平均单步耗时约 49.60 秒。需要指出的是，SFT 与 GRPO 的损失函数定义并不相同，因此不同阶段之间的 loss 数值不适合直接横向比较，更合理的判断依据是各阶段是否稳定完成、权重产物是否完整，以及最终任务级评测结果是否得到改善。

\FloatBarrier
\section{优化效果评测和分析}

\subsection{优化模型任务级评测结果}

在训练完成后，本文将原始 Qwen 3.5 4B 模型作为 Base，将联合 GRPO 后得到的最终权重作为优化模型，在 UAV-DualCog-V1 的 12 个测试场景上进行统一评测。训练集仅来自与测试集场景隔离的图像任务，因此图像任务结果用于衡量优化方法的直接收益，视频任务结果用于观察图像侧证据感知训练是否能够迁移到时间证据推理中。
% 为避免把中间训练阶段误作为最终结论，本节仅围绕 Base 与最终导出的 GRPO 权重模型展开分析；SFT 阶段仅作为训练路线的一部分，不单独进行结果解释。图像任务完整结果按四类任务展开，视频任务完整结果按复合行为、原子行为和地标可见性三类任务展开。

% \autoref{tab:optimization-image-full-results} 和 \autoref{tab:optimization-video-full-results} 中的全局准确率均采用不同任务 Acc 的算术平均计算。这样可以使结果解释直接对应任务级主指标。

\begin{table}[htbp]
    \centering
    \caption{优化模型图像任务实验结果}
    \label{tab:optimization-image-full-results}
    \scriptsize
    \renewcommand{\arraystretch}{0.9}
    \setlength{\tabcolsep}{4pt}
    \resizebox{\textwidth}{!}{%
    \begin{tabular}{l*{13}{c}}
        \toprule
        模型 & 全局准确率 & \multicolumn{6}{c}{自我状态认知} & \multicolumn{6}{c}{环境情况认知} \\
         &  & \multicolumn{3}{c}{地标参照位置推理} & \multicolumn{3}{c}{未来观察预测} & \multicolumn{3}{c}{自我参照位置推理} & \multicolumn{3}{c}{地标驱动动作决策} \\
         &  & Acc & BBox Acc@0.5 & mIoU & Acc & BBox Acc@0.5 & mIoU & Acc & BBox Acc@0.5 & mIoU & Acc & BBox Acc@0.5 & mIoU \\
        \midrule
        Base & 38.99\% & 30.86\% & 12.79\% & 23.14\% & 30.18\% & 11.23\% & 21.69\% & 47.46\% & 6.74\% & 17.71\% & 47.46\% & 6.35\% & 17.96\% \\
        GRPO & 59.77\% & 41.41\% & 74.90\% & 61.91\% & 37.21\% & 66.41\% & 57.36\% & 78.91\% & 52.64\% & 46.23\% & 81.54\% & 50.20\% & 45.63\% \\
        \bottomrule
      \end{tabular}%
    }
\end{table}

\begin{table}[htbp]
    \centering
    \caption{优化模型视频任务实验结果}
    \label{tab:optimization-video-full-results}
    \scriptsize
    \renewcommand{\arraystretch}{0.9}
    \setlength{\tabcolsep}{4pt}
    \resizebox{\textwidth}{!}{%
    \begin{tabular}{l*{10}{c}}
        \toprule
        模型 & 全局准确率 & \multicolumn{3}{c}{复合行为识别} & \multicolumn{3}{c}{原子行为识别} & \multicolumn{3}{c}{地标可见性推理} \\
         &  & Acc/F1 & F1@0.5 & Mean tIoU & Acc/F1 & F1@0.5 & Mean tIoU & Count Acc & F1@0.5 & Mean tIoU \\
        \midrule
        Base & 16.75\% & 6.37\% & 1.31\% & 1.40\% & 10.00\% & 3.04\% & 2.82\% & 33.89\% & 45.18\% & 51.99\% \\
        GRPO & 26.67\% & 24.48\% & 26.90\% & 30.05\% & 24.17\% & 49.08\% & 43.96\% & 31.35\% & 24.23\% & 29.74\% \\
        \bottomrule
      \end{tabular}%
    }
\end{table}

\begin{table}[htbp]
    \centering
    \caption{优化模型在图像与视频任务上的总体增益}
    \label{tab:optimization-overall-results}
    \small
    \begin{tabularx}{\textwidth}{@{}X>{\centering\arraybackslash}p{2.0cm}>{\centering\arraybackslash}p{2.0cm}>{\centering\arraybackslash}p{2.0cm}@{}}
        \toprule
        指标 & Base & GRPO & 提升值 \\
        \midrule
        图像选项准确率 & 38.99\% & 59.77\% & +20.78 pp \\
        图像 BBox Acc@0.5 & 9.27\% & 61.04\% & +51.76 pp \\
        图像 mIoU & 20.12\% & 52.78\% & +32.66 pp \\
        视频全局准确率 & 16.75\% & 26.67\% & +9.91 pp \\
        视频时间 F1@0.5 均值 & 16.51\% & 33.40\% & +16.89 pp \\
        视频 Mean tIoU 均值 & 18.74\% & 34.58\% & +15.84 pp \\
        \bottomrule
    \end{tabularx}
\end{table}

\begin{figure}[htbp]
    \centering
    \includegraphics[width=0.92\textwidth]{fig_optimization_overall_media_gain.png}
    \caption{优化模型在图像与视频核心指标上的总体变化}
    \label{fig:optimization-overall-media-gain}
\end{figure}

\autoref{tab:optimization-image-full-results}、\autoref{tab:optimization-video-full-results}、\autoref{tab:optimization-overall-results} 和 \autoref{fig:optimization-overall-media-gain} 表明，优化模型在图像任务上获得了显著的直接收益：选项准确率从 38.99\% 提升到 59.77\%，说明模型对无人机场景中地标、视角和动作关系的离散判断能力明显增强；BBox Acc@0.5 从 9.27\% 提升到 61.04\%，mIoU 从 20.12\% 提升到 52.78\%，说明优化收益并不只体现在“选项判断正确”，还体现在能够把答案绑定到更可靠的空间证据上。该结果支持了第3章发现的图像侧语义与证据脱节问题。

% 这里需要加强分析，为什么环境可见性计数和区间恢复下降
视频任务也出现了跨媒介迁移收益。视频全局准确率从 16.75\% 提升到 26.67\%，时间 F1@0.5 均值从 16.51\% 提升到 33.40\%，Mean tIoU 均值从 18.74\% 提升到 34.58\%。这说明图像任务上的证据感知训练能够迁移到部分视频时间证据恢复。不过，迁移收益并不均匀：自我行为识别与时间定位显著提升，环境可见性计数和区间恢复下降。因此，本文将视频侧结果解释为“迁移分化”：图像侧训练增强了部分自我行为时间证据能力，但没有同步改善环境可见性时序推理。

\subsection{图像任务直接提升分析}

图像任务是 UAV-DualCog-Train-V1 直接优化的目标媒介。为了观察收益来自哪类能力，本文分别从双认知分支和四类任务两个层面分析优化模型表现。

\begin{table}[htbp]
    \centering
    \caption{优化模型在图像双认知分支上的结果}
    \label{tab:optimization-image-cognition-results}
    \small
    \begin{tabularx}{\textwidth}{@{}p{2.4cm}>{\centering\arraybackslash}p{1.45cm}>{\centering\arraybackslash}p{1.45cm}>{\centering\arraybackslash}p{1.45cm}>{\centering\arraybackslash}p{1.45cm}>{\centering\arraybackslash}p{1.45cm}>{\centering\arraybackslash}p{1.45cm}@{}}
        \toprule
        认知分支 & \multicolumn{2}{c}{选项准确率} & \multicolumn{2}{c}{BBox Acc@0.5} & \multicolumn{2}{c}{mIoU} \\
        \cmidrule(lr){2-3}\cmidrule(lr){4-5}\cmidrule(l){6-7}
        & Base & GRPO & Base & GRPO & Base & GRPO \\
        \midrule
        自我状态认知 & 30.52\% & 39.31\% & 11.99\% & 70.65\% & 22.41\% & 59.64\% \\
        环境情况认知 & 47.46\% & 80.22\% & 6.54\% & 51.42\% & 17.83\% & 45.93\% \\
        \bottomrule
    \end{tabularx}
\end{table}

\autoref{tab:optimization-image-cognition-results} 显示，两条认知分支的提升方式并不相同。环境情况认知的语义选项收益更强，选项准确率从 47.46\% 提升到 80.22\%，说明模型在“地标在哪里”和“应该如何动作”这类外部目标关系判断上最容易吸收训练收益。自我状态认知的选项准确率也有提升，但幅度相对较小，从 30.52\% 到 39.31\%；不过它的空间证据收益更突出，BBox Acc@0.5 从 11.99\% 提升到 70.65\%，mIoU 从 22.41\% 提升到 59.64\%。这说明优化后模型虽然仍然更擅长环境判断，但已经显著改善了自我状态认知中的目标证据定位能力。

\begin{table}[htbp]
    \centering
    \caption{优化模型在四类图像任务上的直接收益}
    \label{tab:optimization-image-task-results}
    \small
    \begin{tabularx}{\textwidth}{@{}p{3.0cm}>{\centering\arraybackslash}p{1.38cm}>{\centering\arraybackslash}p{1.38cm}>{\centering\arraybackslash}p{1.38cm}>{\centering\arraybackslash}p{1.38cm}>{\centering\arraybackslash}p{1.38cm}>{\centering\arraybackslash}p{1.38cm}@{}}
        \toprule
        任务类型 & \multicolumn{2}{c}{选项准确率} & \multicolumn{2}{c}{BBox Acc@0.5} & \multicolumn{2}{c}{mIoU} \\
        \cmidrule(lr){2-3}\cmidrule(lr){4-5}\cmidrule(l){6-7}
        & Base & GRPO & Base & GRPO & Base & GRPO \\
        \midrule
        地标参照位置推理 & 30.86\% & 41.41\% & 12.79\% & 74.90\% & 23.14\% & 61.91\% \\
        未来观察预测 & 30.18\% & 37.21\% & 11.23\% & 66.41\% & 21.69\% & 57.36\% \\
        自我参照位置推理 & 47.46\% & 78.91\% & 6.74\% & 52.64\% & 17.71\% & 46.23\% \\
        地标驱动动作决策 & 47.46\% & 81.54\% & 6.35\% & 50.20\% & 17.96\% & 45.63\% \\
        \bottomrule
    \end{tabularx}
\end{table}

\begin{figure}[htbp]
    \centering
    \includegraphics[width=0.88\textwidth]{fig_optimization_image_task_gain.png}
    \caption{优化模型在四类图像任务上的指标提升}
    \label{fig:optimization-image-task-gain}
\end{figure}

\autoref{tab:optimization-image-task-results} 和 \autoref{fig:optimization-image-task-gain} 进一步说明，优化模型在四类图像任务上都显著提升了空间证据质量。其中，地标参照位置推理的 BBox Acc@0.5 提升 62.11 个百分点，未来观察预测提升 55.18 个百分点，说明优化对自我状态认知中的目标定位尤其有效。环境类任务的选项提升更明显，自我参照位置推理提升 31.45 个百分点，地标驱动动作决策提升 34.08 个百分点，说明模型更容易学习外部地标关系到选项答案之间的映射。综合来看，图像侧优化并不是单一地提高答案准确率，而是把语义判断和空间定位两类证据同时向上拉升，这与 UAV-DualCog 强调“答案必须有证据支撑”的设计目标一致。

\subsection{视频任务迁移提升与分化分析}

由于训练集只包含图像任务，视频任务结果可以用来检验优化模型是否获得跨媒介迁移能力。本文将视频任务拆分为自我行为识别和环境可见性推理两类：前者要求模型理解无人机自身飞行动作并定位对应时间段，后者要求模型统计外部地标可见次数并恢复完整可见区间。

\begin{table}[htbp]
    \centering
    \caption{优化模型在视频任务上的迁移结果}
    \label{tab:optimization-video-transfer-results}
    \small
    \begin{tabularx}{\textwidth}{@{}X>{\centering\arraybackslash}p{1.55cm}>{\centering\arraybackslash}p{1.55cm}>{\centering\arraybackslash}p{1.55cm}>{\centering\arraybackslash}p{1.55cm}>{\centering\arraybackslash}p{1.55cm}>{\centering\arraybackslash}p{1.55cm}@{}}
        \toprule
        任务类型 & \multicolumn{2}{c}{语义/计数主指标} & \multicolumn{2}{c}{时间 F1@0.5} & \multicolumn{2}{c}{时间 IoU} \\
        \cmidrule(lr){2-3}\cmidrule(lr){4-5}\cmidrule(l){6-7}
        & Base & GRPO & Base & GRPO & Base & GRPO \\
        \midrule
        复合自我行为识别 & 6.37\% & 24.48\% & 1.31\% & 26.90\% & 1.40\% & 30.05\% \\
        原子自我行为识别 & 10.00\% & 24.17\% & 3.04\% & 49.08\% & 2.82\% & 43.96\% \\
        环境可见性推理 & 33.89\% & 31.35\% & 45.18\% & 24.23\% & 51.99\% & 29.74\% \\
        \bottomrule
    \end{tabularx}
\end{table}

\begin{figure}[htbp]
    \centering
    \includegraphics[width=0.9\textwidth]{fig_optimization_video_transfer.png}
    \caption{优化模型在视频任务中的迁移分化}
    \label{fig:optimization-video-transfer}
\end{figure}

\autoref{tab:optimization-video-transfer-results} 和 \autoref{fig:optimization-video-transfer} 表明，视频侧迁移具有明显的认知分化。对于自我状态认知，复合自我行为识别的语义主指标从 6.37\% 提升到 24.48\%，时间 F1@0.5 从 1.31\% 提升到 26.90\%；原子自我行为识别的语义主指标从 10.00\% 提升到 24.17\%，时间 F1@0.5 从 3.04\% 提升到 49.08\%。这说明图像侧“目标证据绑定”和结构化输出训练能够迁移到视频中对自身飞行片段的定位，通过加强对环境中动作相关的地标的感知能力来提升对地标相关动作的识别能力，尤其有利于原子行为这类边界更清晰的时序事件。

% 这里可以将可见性的时间区间信息、可见次数做个分析，是不是可见区间累计的占比提升了，但是没有能将可见性有效整合到一段完整的可见区间中，“跟踪”的能力不强，导致可见区间出现中断等情况。
环境情况认知的迁移结果则相反。环境可见性推理的 Count Exact 从 33.89\% 下降到 31.35\%，Segment F1@0.5 从 45.18\% 下降到 24.23\%，时间 IoU 从 51.99\% 下降到 29.74\%。这说明图像任务优化提高了静态目标框证据质量，但没有自然学会“地标在视频中何时出现、持续多久、出现几次”的连续时间建模。换言之，空间证据学习可以迁移到部分自我行为定位，但不能直接替代环境可见性的时序整合能力。

\subsection{跨媒介与双认知综合分析}

\begin{table}[htbp]
    \centering
    \caption{优化模型的跨媒介全局准确率对比}
    \label{tab:optimization-dual-score}
    \small
    \begin{tabularx}{0.78\textwidth}{@{}X>{\centering\arraybackslash}p{2.3cm}>{\centering\arraybackslash}p{2.3cm}>{\centering\arraybackslash}p{2.3cm}@{}}
        \toprule
        模型 & 图像全局准确率 & 视频全局准确率 & 图像-视频差值 \\
        \midrule
        Base & 38.99\% & 16.75\% & 22.24 pp \\
        GRPO & 59.77\% & 26.67\% & 33.10 pp \\
        \bottomrule
    \end{tabularx}
\end{table}

\autoref{tab:optimization-dual-score} 说明，优化模型同时提高了图像全局准确率和视频全局准确率，但两类媒介的提升幅度并不一致。图像全局准确率提升 20.78 个百分点，视频全局准确率提升 9.91 个百分点，图像-视频差值从 22.24 个百分点扩大到 33.10 个百分点。这一结果揭示了能力优化的边界：当前训练路线有效强化了图像空间证据与双认知语义判断，并能向视频行为任务产生一定迁移，但模型能力分布仍更偏向图像媒介，尚未形成跨媒介一致的无人机双认知表示。

从无人机具身推理角度看，这一结果具有明确启发。可靠 UAV 多模态系统不能只在单张图像上判断正确方向或定位地标，也不能只在视频中粗略捕捉目标出现；它需要在自我状态认知和环境情况认知之间建立匹配关系，并在图像空间证据和视频时间证据之间保持稳定转换。优化模型证明了 UAV-DualCog-Train 能显著补强图像侧答案-证据一致性，同时也说明只依赖图像任务训练难以覆盖视频环境可见性推理。因此，后续优化应进一步引入视频任务样本或跨媒介一致性约束，使空间定位收益能够更稳健地迁移到时间区间证据恢复。

% \subsection{本章小结}

% 本章围绕第3章暴露出的结构性问题，系统整理了 UAV-DualCog 的能力优化方法、正式训练过程和最终任务级结果。首先，本文构建了与测试场景严格隔离、在任务类型与双认知分组上保持平衡的 UAV-DualCog-Train-V1 训练集，并将其转换为与 benchmark 对齐的多模态训练记录。其次，本文完成了基于 Qwen 3.5 4B 的多模态 LoRA 监督微调与联合 GRPO 优化方法，明确给出了训练参数、奖励结构、分辨率策略和工程组织方式。再次，本文在 Base 与最终导出的 GRPO 权重模型之间进行任务级对比，结果表明优化模型在图像任务上取得显著直接收益，选项准确率、BBox Acc@0.5 和 mIoU 分别提升 20.78、51.76 和 32.66 个百分点；在视频任务上，按三类任务 Acc 算术平均计算的全局准确率提升 9.91 个百分点，同时呈现迁移分化：自我行为识别与时间定位明显提升，而环境可见性计数与区间定位下降。由此可见，当前能力优化已经有效验证了“证据感知图像训练”对图像双认知能力的补强作用，同时也说明可靠无人机双认知模型仍需要进一步解决图像-视频之间的证据迁移与协同问题。
